\section{Background: Why {\TeX} Limits LLM Workflows}
\label{sec:background}

This section details the frictions summarized in Section~\ref{sec:introduction}.

\subsection{Introduction of {\TeX}}

For decades, {\LaTeX} has been the standard for scientific writing~\cite{lamport1994document}. As an unstructured macro language~\cite{van_der_hoeven_jolly_2020}, however, it presents structural limitations for modern LLM-driven workflows along the axis of machine-tractable document structure \emph{(we defer a brief history of {\TeX} to Appendix~\ref{app:intro})}.

\subsection{Principles of {\TeX} Compilation}

{\TeX} uses a batch-processing model: it reads source code from beginning to end, expanding macros step by step without building an explicit document tree (Abstract Syntax Tree) in memory~\cite{van_der_hoeven_gnu_2001}. Lacking internal state, {\LaTeX} must resolve cross-references and bibliographies through multiple compilation passes, writing temporary data into intermediate files (\texttt{.aux}, \texttt{.toc}, \texttt{.bbl}) and reading them back on the next run \emph{(detailed in Appendix~\ref{app:pr})}. This multi-pass design is inherently slow and opaque, precluding the sub-second, real-time feedback required by modern editors and AI agents~\cite{madje2023typst}.

\subsection{Architectural Defects in {\TeX}}

{\TeX} processes documents as a linear sequence of commands in which structural semantics and typesetting decisions are deeply coupled within a single execution path. Consequently, AI models and static analysis tools cannot reliably extract the document's structure without executing the code \emph{(see Appendix~\ref{app:arch} for a deep analysis)}. This separation problem is fundamental, not incidental: Frank Mittelbach, the technical lead for {\LaTeX}, candidly conceded in an interview~\cite{interview2021} that
\begin{quote}
the separation between local and global in {\LaTeX} seems to be hard. [\ldots] First of all, it is right now hard in {\TeX}\ldots
\end{quote}

The same coupling makes error feedback delayed and ambiguous: the reported location of an error is rarely where the mistake actually happened, and a subtle local defect can irreversibly abort the entire compilation. Listing~\ref{lst:unclosed-frac} shows a typical case---the second argument of \verb|\frac| is missing its closing brace, yet {\TeX} keeps consuming subsequent tokens and only reports a \texttt{Runaway argument?} error at \verb|\end{equation}|, far from the true cause.

\begin{lstlisting}[caption={Unclosed brace: error reported at \texttt{\textbackslash end\{equation\}}, not at its source},label={lst:unclosed-frac}]
\begin{equation}
x=\frac{-b\pm \sqrt{b^{2}-4ac}}{2a
\end{equation}
\end{lstlisting}

Because such logs provide little structural guidance, models are prone to prolonged, often non-terminating debugging loops and hallucinations~\cite{jain2026bibby}. This mechanism constrains LLM-driven writing.

\subsection{UX \& Ecosystem Constraints}

As detailed in Appendix~\ref{app:con}, {\LaTeX}'s ecosystem is too heavy and fragmented: a full TeX Live distribution exceeds 6GB, hindering deployment in cloud environments or sandboxes, while different engines (pdfLaTeX, XeLaTeX, LuaLaTeX) handle fonts and Unicode differently, yielding inconsistent results. While {\LaTeX} retains clear technical strengths in typographic quality and mathematical expressivity, its continued dominance is also sustained by established user habits and publisher requirements~\cite{van_der_hoeven_jolly_2020}; for LLM-driven workflows specifically, these factors leave its structural limitations unaddressed. This motivates a more scalable, LLM-friendly data representation.